\section{Механика}
След като дефинирахме основно приближение в нашия, можем да преминем към изграждането на математическия
модел и последващото му решаване. Такъв един модел се поддава добре на разглеждане от гледна точка на
диференциалната геометрия. Физичният смисъл на това е, че се фокусираме върху взаимодействието на материални
точки (,,частици``) в рамките на апроксимацията на непрекъснатите среди (следствията от това приближение
ще бъдат дискутирани по-нататък в текста). 
Освен това ще бъдат разгледани само дисперсионните взаимодействия (т.нар $r^{-6}$ потенциал).

\begin{figure}[h]
    \centering 
    \subfigure[Схема на взаимодействието]{\label{fig:cylinder_front}\includegraphics[width=0.48\linewidth]{cyl_fig_front.png}}
    \subfigure[Ортогонална проекция]{\label{fig:cylinder_top}\includegraphics[width=0.48\linewidth]{cyl_fig_top.png}}
    \caption{
        Геометрична схема на разглежданата система. $r_{ij}$ - растояние между вътрешната \textbf{\textit{i}} и външната \textbf{\textit{j}} ,,частици``. $z_{ij}$ и $\rho_{ij}$ са съотвените проекции. $r_{ij}^2 = \rho_{ij}^2 + z_{ij}^2$. \textbf{\textit{r}} - радиус на цилиндъра.}
    \label{fig:cylinder_schematic}
\end{figure}

Такова разглеждания са правени много за случая на плоскопарелелните течни филми \cite{rusanov_shchekin}, като най-интересни са ,,тънките`` филми поради значителното влияние на капилярните
явления в тези системи. ,,Тънки`` (капилярни) са тези системи, чиито размери поне по едно от направленията са достатъчно малки, 
за да се наблюдава припокриване между двете междуфазови преходни области. В тях няма изотропен/обемен флуид, т.е. свойствата в тях са изцяло анизотропни. 
Всички ,,интересни`` капилярни явления всъщност са следствие от тази анизотропност на средата. 

Тук следва да отблежим, че механичния подход за тези системи е особено подходящ, тъй като в термодинамичния такъв губим детайлната информация за анизотропията - в случая на плоските филми
,,заместваме`` системата с опростена такава - обемен флуид ,,затворен`` между две математически равнини. От друга страна термодинамичния модел е удобен от практическа гледна точка - параметрите и енергиите
са адитивни. По тази причина оснвният ни анализ ще бъде механичен, а термодинамичните величини - напрежение на ,,филма``, повърхностно напрежение, разклинящо налягане и т.н. ще бъдат изведени като следствия
от получената по механичния подход енергия на взаимодействие.

\subsection{Получаване на енергията на взаимодействие}
Основна стъпка към анализа на система е получаването на плътността на енергията на взаимодействие между частиците. 
Дисперсионните взаимодействия имат адитивен характер \cite{israelachvili}, което ни позволява да представим нашата система като
,,потопена`` в един резервоар (в случая с безкрайни размери) или:
\begin{equation}
    w_{\infty}  = w + \tilde{w}
\end{equation}
Където $\tilde{w}$ е плътността на взаимодействието на елемент от разглежданата система \textbf{\textit{i}} ($0 < \rho_{i} < r$) (оградена от плътния контур на \autoref{fig:cylinder_schematic}) с елемент \textbf{\textit{j}} от съседната на цилиндъра част от безкрайна система ($r < \rho_{j} <  \infty$).
Безкрайния резервоар е по дефиниция хомогенен, т.е. неговата плътност на енергията е постоянна. Това позволява да заместваме съвсем просто енергията на цилиндъра с енергията на допълнителната система в уравненията, без загуба на обобщеност.
Още повече това ни позволява да останем в рамките на модела на непрекъснатите среди като избегнем работата с взаимодействия на къси разстояния, за които въпросния модел не е валиден, тъй като дискретната структура на системата започва да има силно отношение върху разглежданията.
В моделите, за да се избегнат сингулярности, които ще бъдат изведени са въвежда ,,най-малко разстояние`` $a$ (или нормирано към радиуса на цилиндъра $a_{c}$),
което стандартно се интерпретира като ,,молекулен диаметър``, макар че придаването на този физичен смисъл става \textit{a posteriori} \cite{israelachvili} и следва да се внимава 
при интерпретацията на резултатите, когато се работи с него.